\scanpage{15}
\mspage{13}
\uline{Dém} a) Comme \(\pi'_0\unsure{\longrightarrow|\mathrel{-}}\pi_0\) est évidemment un isom (car \(k'/k\) est l'ext. triviale) il suffit de prouver que \(I'\xrightarrow{\sim}I\), et comme \(I'_t\xrightarrow{\sim}I_t\) (car \(e=1\)) il suffit de prouver que \(P'\xrightarrow{\sim}P\).

b) On sait que \(\pi'\mathrel{-}\pi\) est surjectif, i.e pour tte ext. finie sép. \(K_1/K\), \(K_1\otimes_K\widehat K\) est un \uline{corps} : en effet, la clôture normale \(V_1\) de \(V\) dans \(K_1\) est \uline{finie}, donc \(V_1\otimes_V\widehat V\simeq\widehat V_1\), \deleted{\unsure{est|et}} \inserted{c'est donc} un anneau de val. discrète (\(V_1\) en étant un), donc \(K_1\otimes_K\widehat K=\widehat V_1\otimes_V\widehat K\) est un corps.

\deleted{\unsure{D'autre|Du reste|Dans} Quitte à remplacer \(V,V'\) par \(\widetilde V,\widetilde V'\) et à comparer aussi leurs complétés \illegible{2} \(\widehat{\widetilde V}\simeq\widehat{\widetilde V'}\) de ces derniers, \unsure{on est ramené|on est amené}}

c) On \deleted{\unsure{en|est}} conclut \inserted{de a) et b)}, sur les \inserted{\unsure{hom de|hom. de}} suites exactes, que \(I'\mathrel{-}I\) est surjectif, puis \(P'\mathrel{-}P\) \uline{surjectif}. Il reste à prouver que \(P'\mathrel{-}P\) est injectif, et pour ceci que \(\pi'\mathrel{-}\pi\) est injectif, \deleted{Or} i.e. que toute extension séparable finie \(K'_1\) de \(K'=\widehat K\) provient d'une ext. sép. finie de \(K\). Or ceci résulte du fait connu (``continuité des racines p.r. aux coefficients'') que deux équations séparables \deleted{\unsure{de|dans}} à coefficients dans \(\widehat K\) \uline{voisines} définissent des \deleted{\unsure{corps|extensions}} \(\widehat K\)-Algèbres isomorphes.

% archive sheet number: 14
